\section{Introducción}

El continuo avance de las arquitecturas de hardware ha impulsado la evolución de la Computación de Alto Rendimiento (HPC), permitiendo resolver problemas matemáticos de dimensiones masivas en tiempos computacionalmente viables. Un exponente fundamental de estos desafíos es la multiplicación de matrices, una operación algebraica cuya complejidad de \(\mathcal{O}(n^3)\) impone una carga de procesamiento extrema a medida que crecen las dimensiones de los datos. A pesar de la aparente simplicidad de su implementación iterativa básica, la optimización de este algoritmo es un área de constante investigación, dado que constituye el núcleo estructural de aplicaciones modernas críticas, tales como el álgebra lineal computacional, la simulación física y el entrenamiento de modelos de inteligencia artificial \parencite{Dongarra1990BLAS}.

Históricamente, el aumento de la frecuencia de reloj en los microprocesadores permitía acelerar la ejecución de algoritmos secuenciales de forma directa. Sin embargo, tras alcanzar los límites físicos y térmicos del silicio, la industria se desplazó hacia arquitecturas multinúcleo que exigen un paradigma de programación paralela \parencite{Patterson2017Computer}. Bajo este enfoque, el diseño de software eficiente requiere distribuir la carga de trabajo de manera concurrente. Para lograrlo en el lenguaje C, los sistemas operativos ofrecen distintos mecanismos, destacando la paralelización ligera mediante hilos (compartiendo el mismo espacio de memoria) y la paralelización pesada mediante la creación de procesos independientes; cada uno con implicaciones directas sobre el rendimiento global y el consumo de recursos \parencite{Silberschatz2018Operating}.

El presente trabajo tiene como objetivo principal analizar, comparar y documentar el desempeño del algoritmo de multiplicación de matrices bajo tres enfoques de programación: una implementación secuencial, una concurrente basada en hilos mediante la librería Pthreads, y una paralela basada en la duplicación de procesos mediante la llamada al sistema \texttt{fork}. A través de experimentación empírica variando el tamaño de las matrices y la cantidad de unidades de procesamiento asignadas, se medirán los tiempos reales de ejecución. Finalmente, apoyándose en métricas estandarizadas como el \textit{Speedup} (aceleración) \parencite{Hager2010HPC}, se busca comprobar matemáticamente la eficacia de la paralelización para reducir los tiempos de cómputo, así como evidenciar los cuellos de botella generados por la gestión de memoria y el límite de los recursos físicos del hardware.
